%! Adding \& Subtracting Rational Expressions — Unit 4, Lesson 4.3 — Exit Ticket (Answer Key)
\documentclass[10pt]{article}
\usepackage{algebra2-article}
\usepackage{algebra2-key}   % pulls in -boxes; do NOT also load -boxes
\newcommand{\TallMath}[1]{$\displaystyle #1\rule[-1.4em]{0pt}{3.2em}$}

\begin{document}

\pageheader{Unit 4, Lesson 4.3}{Exit Ticket --- Answer Key}
\namedateperiod

\vspace{0.08in}

No notes. Show the LCD and the combined numerator for question 1.

\vspace{0.06in}

\begin{enumerate}[leftmargin=1.8em, itemsep=9pt]
  \item \textbf{Combine and restrict.} \; $\dfrac{5}{x-2}-\dfrac{20}{x^2-4}$

        \vspace{4pt}
        Denominators, factored: \ans{$(x-2)$ and $(x-2)(x+2)$} \quad LCD: \ans{$(x-2)(x+2)$}

        \vspace{4pt}
        Combined numerator, factored: \ans{$5x-10=5(x-2)$} \quad Simplest form: \ans{$\frac{5}{x+2}$}

        \vspace{4pt}
        All restrictions: $x\neq$ \ans{$2,\ -2$}. \; One of them is a \textbf{hole} and one is a
        \textbf{wall}. \\[3pt]
        Hole at $x=$ \ans{2}, height \ans{$\frac54$} \qquad Wall at $x=$ \ans{$-2$}

  \item \textbf{Explain.} A classmate wrote
        $\dfrac{4x+1}{x-6}-\dfrac{x-5}{x-6}=\dfrac{4x+1-x-5}{x-6}=\dfrac{3x-4}{x-6}$.
        Find the error, give the correct answer with its restriction, and prove your version is the
        right one using $x=0$.
        \ansline{The second numerator needed parentheses: $-(x-5)=-x+5$, not $-x-5$. Correct:}
        \ansline{$\frac{(4x+1)-(x-5)}{x-6}=\frac{3x+6}{x-6}=\frac{3(x+2)}{x-6}$, with $x\neq6$.}
        \ansline{At $x=0$ the original is $\frac{1}{-6}-\frac{-5}{-6}=-\frac16-\frac56=-1$; mine gives $\frac{6}{-6}=-1$; the classmate's gives $\frac23$.}

  \item \textbf{SOL-style multiple choice.} Which is the simplest form of
        \; \TallMath{\frac{\;\dfrac{1}{x}-\dfrac{1}{5}\;}{\;x-5\;}} \; with \emph{all} of its
        restrictions?
        \begin{enumerate}[label=\Alph*., leftmargin=1.9em, itemsep=1pt]
          \item $\dfrac{1}{5x}$, \; $x\neq0,\,5$
          \item $-\dfrac{1}{5x}$, \; $x\neq0$
          \item $-\dfrac{1}{5x}$, \; $x\neq0,\,5$
                \quad \textcolor{keyred}{\textbf{$\leftarrow$ correct}}
          \item $\dfrac{5-x}{5x-25}$, \; $x\neq5$
        \end{enumerate}
        \vspace{2pt}
        {\small The top combines to $\dfrac{5-x}{5x}$; dividing by $x-5$ gives
        $\dfrac{5-x}{5x(x-5)}=\dfrac{-(x-5)}{5x(x-5)}=-\dfrac{1}{5x}$, with $x\neq0$ from the inner
        denominator and $x\neq5$ from the main fraction bar. \textbf{A} loses the $-1$ from the
        opposite binomials, \textbf{B} forgets the main bar's restriction, and \textbf{D} never
        simplified and then read its restriction off its own denominator.}
\end{enumerate}

\begin{teachernote}
Graded for completion --- ``mistakes happen, blanks don't.'' Sort into four piles: \textbf{(1) got it};
\textbf{(2) sign error} --- item 2 wrong, or item 1's numerator given as $5x+10$; \textbf{(3) LCD built
as a product} --- item 1 done over $(x-2)(x^2-4)$, which is not wrong but leaves a mess and usually
loses the cancellation; \textbf{(4) restrictions incomplete} --- chose \textbf{B}, or gave item 1 only
$x\neq-2$ (read off the answer, which is Lesson 4.1's error at its third appearance).

\vspace{3pt}
Pile 2 is the one to re-teach: the parenthesis is the single most common lost point on this standard,
and it does not go away on its own. Pile 4 decides nothing about tomorrow --- 4.4 is graphing the parent
$y=1/x$ --- but it decides a great deal about \textbf{4.6}, where a solution landing on any excluded
value is exactly what ``extraneous'' means. Note who is in pile 4 and check them again then.

\vspace{3pt}
Item 1's hole/wall question is the graph-reading half of the item and should not be treated as a bonus:
students who can combine but cannot say which excluded value survives into the answer have the algebra
without the function.
\end{teachernote}

\end{document}
